% Brief: the setup of the lesson next to it, as steps for another person's agent.
% Compile with:  latexmk -jobname=working-with-claude-brief brief.tex

\documentclass[11pt,a4paper]{article}

\usepackage{fontspec}
\usepackage[margin=2.2cm]{geometry}
\usepackage{xcolor}
\usepackage[most]{tcolorbox}
\usepackage{listings}
\usepackage{twemojis}
\usepackage{alfdoc}
\usepackage{enumitem}
\usepackage{booktabs}
\usepackage{microtype}  % font expansion, so lines full of unbreakable paths keep even spacing
\usepackage[colorlinks=true]{hyperref}

\setlength{\parindent}{0pt}
\setlength{\parskip}{0.5em}
\setlength{\emergencystretch}{3em}
\newcommand{\home}{\textasciitilde}
\renewcommand{\slash}{/\allowbreak}  % paths break at their slashes

% Review: “Not just Claude (any agent work, but yes, Claude is my default agent).”
\title{\emo{file folder}\ Tidying a setup with an agent\\[0.3em]
  \large A brief for your agent, from any starting point to one home per project}
% Review: “Let's keep the written by X for Y, where Y is my student.”
\author{\here[for whom]Brief for my student, written by Claude}
\date{Thu 12 Nov 2026}

\begin{document}
\maketitle
\vspace{-1.5em}

\begin{changes}
  \item It says who wrote it and for whom.
  \item The brief is for any agent, and Claude is mine.
  \item Paths break at their slashes.
  \item What a project folder holds: its note and links to its files.
  \item \name{To-delete/} is one word.
  \item The inbox is whatever folder the owner saves into.
  \item The knowledge base, and links that gather scattered places.
\end{changes}

% Review: “Not just Claude (any agent work, but yes, Claude is my default agent).”
\here[any agent]This brief is for an AI agent, such as Claude Code or Codex, working with someone whose files and habits make things hard to find. Claude is my default agent. The starting point can be any of the usual ones: everything saved in \name{Downloads/}, a desktop covered in files, one huge folder with no structure, the same project spread across \name{Dropbox}, \name{Google Drive}, and email, notes in several apps or none, or every agent session started from the home folder, so that all projects share one memory. My lesson \emph{How I work with Claude}~\cite{wwc} describes the setup I ended up with: one home per project, rules in one place, and a notes vault for what spans projects. That lesson is the shape to aim for. This brief is the way to get there, one step at a time, with the owner deciding at every step.

\subsection*{1.\ Look before touching anything}

Find where the person actually keeps things: the folders they save into (often \name{Downloads/}, \name{Desktop/}, and \name{Documents/}), the cloud services they use, and where their notes and repositories live. Survey those places without changing them. Group the files by kind (papers, slides, code, datasets, installers, images, archives), by age, and by size, and find exact duplicates by their checksum. Show the owner a short summary: how many files of each kind and where, the largest ones, the duplicates, and the installers (\texttt{.dmg}, \texttt{.pkg}) that have already done their job.

\subsection*{2.\ Ask what the work is}

Ask the owner what they are working on now (courses, research projects, applications), what they look after over time (teaching, health, finances), and what they keep only as reference. These answers become the folders, so the layout follows their work and the files are sorted into it.

\subsection*{3.\ Propose a layout and wait for an answer}

% Review: “Sec 3 has bad spacing (needs breakable file paths).”
\here[breakable paths]Propose a small layout following PARA~\cite{para}, Tiago Forte's Projects, Areas, Resources, and Archive: \name{\home/PARA\slash Projects\slash <name>/} for work with an end, \name{\home/PARA\slash Areas/} for what continues, \name{\home/PARA\slash Resources/} for reference, and \name{\home/PARA\slash Archive/} for what is finished. Git repositories go in one folder outside any cloud-synced folder, such as \name{\home/Work\slash GitHub/}, and temporary files go in \name{\home/Work\slash Scratch/}. Show the layout as a short nested list, and create it only after the owner agrees.

% Review: “I'm not sure it explains what Projects contains.”
% Review: “There is no mention of md-Wiki nor the file tree of links. That's how I go around having
% things that were and are scattered in multiple locations under the same header.”
\here[a project's folder]A project's files often live in several places at once: a repository, a data folder, a shared cloud folder, and a note. Each folder in \name{Projects/} gathers them under one name as a folder of links. In mine, \name{repo} links to the repository in \name{\home/Work\slash GitHub/}, \name{local} to the data folder, \name{g-drive} to the shared Google Drive folder, and \name{note.md} to the project's note, so opening the project's folder shows everything it uses, wherever it is stored. Links are one way to do this, and the owner may instead move the files themselves into the project's folder.

\subsection*{4.\ Move files in small batches}

% Review: “No spaces in To delete. Also, usually Downloads is not what we opened with.”
Move one group at a time, for example all the papers, after showing the owner which file goes where. Write every move to a log, \name{\home/PARA\slash moves.tsv}, with the old path and the new one, so any batch can be undone. Put duplicates and spent installers in \here[no space]\name{\home/PARA\slash To-delete/}, where the owner empties them after a look. \here[the owner's folder]The folder the owner habitually saves into becomes an inbox: anything new lands there and gets sorted the same way.

\subsection*{5.\ Give each project its own home}

Start each agent session from the project's folder, so that the agent's memory for that project stays with it, and keep a short \name{note.md} there with the goal, the next action, and a dated log. \here[knowledge base]The notes themselves live together in one knowledge base, a folder of Markdown notes linked to each other; mine is called \name{md-Wiki/}. Each project's \name{note.md} is a link to its note there, so the knowledge base holds what spans projects, such as people, tools, and ideas, next to every project's note. Write the owner's global rules for the agent last, once a few weeks of work show how they actually like to work: for Claude Code they go in \name{\home/.claude\slash CLAUDE.md}, for Codex in \name{\home/.codex\slash AGENTS.md}.

\begin{thebibliography}{PARA}
\bibitem[WwC]{wwc} Alf, \textit{How I work with Claude}.
  \url{https://github.com/Atcold/learning-with-claude/tree/master/working-with-claude}
\bibitem[PARA]{para} Tiago Forte, \textit{The PARA method}, Forte Labs.
  \url{https://fortelabs.com/blog/para/}
\end{thebibliography}

\revisiontable

\end{document}
